
% ════════════════════════════════════════════
% 模型建立（5.1.3.模型建立.tex，三级子块）
% ════════════════════════════════════════════
\subsubsection{单日购电—储能调度的线性规划模型}

针对问题一，本节从能量平衡与储能状态递推出发，建立微网日购电计划的线性规划模型。

\textbf{步骤1：决策变量与目标函数}

设第 $k$ 个间隔的购电功率为 $g_k$，储能充电功率为 $c_k$，放电功率为 $d_k$，储电量为 $E_k$，其中 $k=0,1,\dots,143$，$E_k$ 为第 $k$ 个间隔起始时刻的储电量。以微网全天购电费用为目标函数，即各间隔购电能量与电价乘积之和：
\begin{equation}
    \min \; Z=\sum_{k=0}^{143} g_k\,p_k\,\Delta t
\end{equation}
其中 $g_k\geq 0$ 为购电功率，$p_k$ 与 $\Delta t$ 分别为第 $k$ 个间隔的交易电价与时间步长。

\textbf{步骤2：功率平衡约束}

由功率守恒，每个间隔内购电功率、光伏功率与储能放电功率之和须不低于负载功率与储能充电功率之和，即允许在光伏过剩时弃光：
\begin{equation}
    g_k + S_k + d_k - c_k \geq L_k,\quad k=0,1,\dots,143
\end{equation}
其中 $S_k$、$L_k$ 分别为光伏预测功率与小区负载功率。光伏发电边际成本为零，最优解会自动优先消纳光伏，仅在光伏与储能放电不足时才从外网补充购电。

\textbf{步骤3：储能状态递推与运行约束}

储能荷电状态随充放电逐间隔更新。充、放电效率均为 $\eta$，充电时只有 $\eta$ 倍能量存入电池，放电时向负荷释放单位能量需从电池取出 $1/\eta$ 倍，故递推关系为
\begin{equation}
    E_{k+1}=E_k+\eta\,c_k\,\Delta t-\frac{1}{\eta}\,d_k\,\Delta t,\quad k=0,1,\dots,143
\end{equation}
储能还须满足容量上下限与变换器额定功率限制
\begin{equation}
    E_{\min}\leq E_k\leq E_{\max},\quad 0\leq c_k,d_k\leq P_{\max}
\end{equation}
其中 $\eta=0.9$，$E_{\min}=1200$、$E_{\max}=10800$（kWh），$P_{\max}=5000$ kW。

\textbf{步骤4：储电初值与日循环约束}

附录1给定 2025 年 1 月 1 日 0:00 的储电量为 6000 kWh，且要求 0:00 与 24:00 储电量相同，以保证相邻日策略的可重复性，故
\begin{equation}
    E_0=6000,\quad E_0=E_{144}
\end{equation}
该模型共有 $4\times 144+1=577$ 个连续变量，其中 $E_0$ 由附录1固定为 6000 kWh；全部约束均为线性等式或不等式，属于标准的线性规划问题\cite{dantzig1963}，可用单纯形法或内点法精确求解，所得解即全局最优解。
